\documentclass[11pt,a4paper,fontset=fandol]{ctexart}

\usepackage[a4paper,top=2.25cm,bottom=2.45cm,left=2.25cm,right=2.25cm,headheight=18pt,headsep=0.55cm,footskip=24pt]{geometry}
\usepackage{amsmath,amssymb,amsthm}
\usepackage{fancyhdr}
\usepackage{lastpage}
\usepackage{enumitem}
\usepackage{titlesec}
\usepackage{xcolor}
\usepackage{hyperref}
\usepackage{bookmark}
\usepackage[most]{tcolorbox}
\usepackage{setspace}
\usepackage{array}

\hypersetup{
  colorlinks=true,
  linkcolor=blue!50!black,
  urlcolor=blue!50!black
}

\setstretch{1.13}
\setlength{\parindent}{2em}
\setlength{\parskip}{0.25em}
\allowdisplaybreaks
\setlist[enumerate]{itemsep=0.45em, topsep=0.35em, leftmargin=2.4em}
\setlist[itemize]{itemsep=0.25em, topsep=0.25em, leftmargin=1.9em}

\definecolor{mainblue}{RGB}{36,83,140}
\definecolor{lightblue}{RGB}{241,246,252}
\definecolor{lightgray}{RGB}{246,246,246}
\definecolor{rulegray}{RGB}{110,118,128}

\titleformat{\section}
  {\Large\bfseries\color{mainblue}}
  {\thesection.}
  {0.45em}
  {}
  [\vspace{0.2em}\color{rulegray}\titlerule]
\titlespacing*{\section}{0pt}{1.1em}{0.7em}

\pagestyle{fancy}
\fancyhf{}
\fancyhead[L]{\small 容斥原理周测 A 卷}
\fancyhead[C]{\small 组合专题：基础计数与坏事件模型}
\fancyhead[R]{\small 刘文博}
\fancyfoot[C]{\small 第 \thepage 页 / 共 \pageref{LastPage} 页}
\renewcommand{\headrulewidth}{0.55pt}
\renewcommand{\footrulewidth}{0.35pt}

\newtcolorbox{infobox}[1]{
  breakable,
  colback=lightblue,
  colframe=mainblue,
  boxrule=0.7pt,
  arc=1mm,
  title=\textbf{#1},
  fonttitle=\bfseries
}

\newenvironment{solution}{\par\noindent\textbf{参考解答：}}{\hfill$\square$\par}
\newcommand{\answerspace}{\vspace{2.3cm}}
\newcommand{\biganswerspace}{\vspace{3.1cm}}
\newcommand{\Dn}{D_n}

\begin{document}

\thispagestyle{empty}
\begin{center}
{\fontsize{22}{28}\selectfont\bfseries 容斥原理周测 A 卷\par}
\vspace{0.4cm}
{\large 适合小学高年级至初中低年级数学竞赛训练\par}
\end{center}

\begin{infobox}{考试说明}
\begin{itemize}
  \item 测试时间：70--75 分钟
  \item 满分：100 分
  \item A 卷侧重两集合、三集合、补集法和基础坏事件模型，建议先写清楚“设了哪些集合/事件”。
\end{itemize}
\end{infobox}

\section{试题}

\begin{enumerate}
  \item （10 分）在 $1$ 到 $100$ 中，是 2 的倍数或 3 的倍数的整数有多少个？
  \answerspace

  \item （12 分）在 $1$ 到 $120$ 中，是 3 的倍数或 5 的倍数或 8 的倍数的整数有多少个？
  \answerspace

  \item （12 分）在 $1$ 到 $200$ 中，不是 2 的倍数也不是 3 的倍数的整数有多少个？
  \answerspace

  \item （12 分）5 个人甲、乙、丙、丁、戊排成一行。要求甲不站在第 1 个位置，乙不站在第 2 个位置。共有多少种排法？
  \answerspace

  \item （12 分）6 个人甲、乙、丙、丁、戊、己排座位。要求甲不坐 1 号位，乙不坐 2 号位，丙不坐 3 号位，丁不坐 4 号位，另外两人没有限制。共有多少种排法？
  \biganswerspace

  \item （14 分）6 个不同数字 $1,2,3,4,5,6$ 排成一行，要求 $1,2,3$ 都不在自己的位置上，而 $4,5,6$ 没有限制。共有多少种排法？
  \biganswerspace

  \item （14 分）用容斥原理计算 $D_5$。
  \answerspace

  \item （14 分）8 个不同数字排成一行，至少有 1 个数字在原位，共有多少种排法？
  \biganswerspace
\end{enumerate}

\newpage
\section{详细解答}

\begin{enumerate}
  \item \begin{solution}
  设
  \[
  A=\{\text{2 的倍数}\},\qquad B=\{\text{3 的倍数}\}.
  \]

  则
  \[
  |A|=50,
  \qquad
  |B|=33,
  \qquad
  |A\cap B|=\left\lfloor\frac{100}{6}\right\rfloor=16.
  \]

  所以
  \[
  |A\cup B|=50+33-16=67.
  \]
  \end{solution}

  \item \begin{solution}
  设 $A,B,C$ 分别表示“3 的倍数”“5 的倍数”“8 的倍数”。

  单个集合：
  \[
  |A|=40,
  \qquad
  |B|=24,
  \qquad
  |C|=15.
  \]

  两两交集：
  \[
  |A\cap B|=\left\lfloor\frac{120}{15}\right\rfloor=8,
  \]
  \[
  |A\cap C|=\left\lfloor\frac{120}{24}\right\rfloor=5,
  \qquad
  |B\cap C|=\left\lfloor\frac{120}{40}\right\rfloor=3.
  \]

  三者交集：
  \[
  |A\cap B\cap C|=\left\lfloor\frac{120}{120}\right\rfloor=1.
  \]

  所以
  \[
  40+24+15-8-5-3+1=64.
  \]

  因此共有
  \[
  64
  \]
  个整数。
  \end{solution}

  \item \begin{solution}
  先设
  \[
  A=\{\text{2 的倍数}\},\qquad B=\{\text{3 的倍数}\}.
  \]

  那么
  \[
  |A|=100,
  \qquad
  |B|=66,
  \qquad
  |A\cap B|=33.
  \]

  所以 1 到 200 中，是 2 的倍数或 3 的倍数的整数有
  \[
  100+66-33=133
  \]
  个。

  总共有 200 个整数，因此既不是 2 的倍数也不是 3 的倍数的有
  \[
  200-133=67.
  \]
  \end{solution}

  \item \begin{solution}
  设
  \[
  A=\text{“甲站第 1 个位置”},\qquad B=\text{“乙站第 2 个位置”}.
  \]

  总排法有
  \[
  5!=120
  \]
  种。

  单个坏事件：
  \[
  |A|=|B|=4!=24.
  \]

  两个坏事件同时发生：
  \[
  |A\cap B|=3!=6.
  \]

  因此满足要求的排法为
  \[
  120-24-24+6=78.
  \]
  \end{solution}

  \item \begin{solution}
  设
  \[
  A_1=\text{“甲坐 1 号位”},
  A_2=\text{“乙坐 2 号位”},
  A_3=\text{“丙坐 3 号位”},
  A_4=\text{“丁坐 4 号位”}.
  \]

  总排法有
  \[
  6!=720
  \]
  种。

  单个坏事件共有
  \[
  \binom41 5!=4\times 120=480
  \]
  种。

  两个坏事件同时发生共有
  \[
  \binom42 4!=6\times 24=144
  \]
  种。

  三个坏事件同时发生共有
  \[
  \binom43 3!=4\times 6=24
  \]
  种。

  四个坏事件同时发生共有
  \[
  \binom44 2!=2.
  \]

  因此所求为
  \[
  720-480+144-24+2=362.
  \]
  \end{solution}

  \item \begin{solution}
  设
  \[
  A_1=\text{“1 在第 1 位”},\quad
  A_2=\text{“2 在第 2 位”},\quad
  A_3=\text{“3 在第 3 位”}.
  \]

  总排法有
  \[
  6!=720
  \]
  种。

  单个坏事件共有
  \[
  3\times 5!=360
  \]
  种。

  两个坏事件同时发生共有
  \[
  \binom32 4!=3\times 24=72
  \]
  种。

  三个坏事件同时发生共有
  \[
  3!=6
  \]
  种。

  所以满足要求的排法数为
  \[
  720-360+72-6=426.
  \]
  \end{solution}

  \item \begin{solution}
  设 $A_i$ 表示“第 $i$ 个对象在原位”，其中 $i=1,2,3,4,5$。

  那么
  \[
  D_5=5!-\binom51 4!+\binom52 3!-\binom53 2!+\binom54 1!-\binom55 0!.
  \]

  代入得
  \[
  D_5=120-120+60-20+5-1=44.
  \]
  \end{solution}

  \item \begin{solution}
  “至少有 1 个数字在原位”最适合用补集法。

  总排法有
  \[
  8!=40320
  \]
  种。

  一个都不在原位的排法数是
  \[
  D_8.
  \]
  由递推公式：
  \[
  D_7=1854,
  \qquad
  D_8=7(D_7+D_6)=7(1854+265)=14833.
  \]

  因此所求为
  \[
  40320-14833=25487.
  \]
  \end{solution}
\end{enumerate}

\end{document}
